%% ======================================================================
%%  Bibliography (pp. 106-107)
%% ======================================================================
\clearpage
\folio{106}%
\global\let\currentseccode\empty
\global\let\pagefirstsec\empty
\renewcommand{\chaptertitle}{}%
\def\currentsectitle{}%
\thispagestyle{emptyhead}%
\phantomsection
\pdfbookmark[0]{Bibliography}{fm:bibliography}%
\label{fm:bibliography}%
\addtocontents{toc}{\string\fml{BIBLIOGRAPHY}{\thepage}{fm:bibliography}}%
\begin{center}{\Large Bibliography}\end{center}%
\vspace{0.8\baselineskip}%
\noindent\footnotesize*\ \ The references preceded by a number between
brackets are referred to in the text.%
\vspace{1.0\baselineskip}%
\par
\normalsize

\begingroup
\setlength{\parindent}{0pt}%
\setlength{\parskip}{0.45\baselineskip}%

E.\thinspace ARTIN, \emph{Theory of algebraic numbers} (G.\ Striker,
Schildweg 12, Göttingen, Germany---1957) (gives a presentation of valuation
theory, very elegant; numerous examples).

H.\thinspace HASSE, \emph{Zahlentheorie} (Akademische Verlag, Berlin, 1949)
(a large and very complete work).

H.\thinspace HASSE, \emph{Vorlesungen über Zahlentheorie} (Springer, 1964)
(describes many aspects of number theory).

G.\thinspace H.\thinspace HARDY and E.\thinspace M.\thinspace WRIGHT,
\emph{An introduction to the theory of numbers} (Clarendon Press, Oxford,
1965) (profound and beautiful; shows a remarkable aesthetic sense in the
choice of material).

E.\thinspace HECKE, \emph{Vorlesungen über die Theorie der algebraischen
Zahlen} (Chelsea, New York, 1948) (a classic; elegant and complete).

S.\thinspace LANG, \emph{Algebraic numbers} (Addison-Wesley, 1964) (a little
book---very dense and concentrated).

S.\thinspace LANG, \emph{Diophantine geometry} (Interscience Tract no.~11,
J.\ Wiley, New York, 1962) (oriented toward diophantine equations; gives a
nice indication of their connection with algebraic geometry).

O'MEARA, \emph{Introduction to quadratic forms} (Springer, 1963) (an
efficacious exposition of algebraic number theory, followed by one of its
most beautiful applications and motivations).

J.\thinspace P.\thinspace SERRE, \emph{Corps locaux} (Hermann, Paris, 1962)
(the emphasis here is on $p$-adic fields; a very clear and lucid exposition
of the most recent algebraic methods in number theory; very rich in content;
numerous examples).

E.\thinspace ARTIN and J.\thinspace TATE, \emph{Class-field theory}
(W.\ A.\ Benjamin, 1968) (the most modern exposition of the famous
class-field theory, i.e.\ the theory of abelian extensions of number
fields).

A.\thinspace WEIL, \emph{Basic number theory} (Springer, 1967) (utilizes the
fruitful method of adèles and treats simultaneously number fields and
function fields).

Z.\thinspace I.\thinspace BOREVICH and I.\thinspace R.\thinspace SHAFAREVICH,
\emph{Number theory} (Academic Press, 1966) (very complete; excellent
chapters on complex and $p$-adic analytic methods; numerous numerical
tables).

\endgroup
\vspace{1.2\baselineskip}

\begin{list}{}{\leftmargin=2.6em\labelwidth=2.2em\labelsep=0.4em
  \itemindent=0pt\parindent=0pt\parsep=0.45\baselineskip\topsep=0pt}
\item[\textbf{[1]}] \phantomsection\label{bib:1}N.\thinspace BOURBAKI, \emph{Algèbre} (Paris, Hermann)
(especially Chapter V which discusses fields, Chapter VI which discusses
divisibility, and Chapter VII which discusses modules over principal ideal
rings).
\item[\textbf{[2]}] \phantomsection\label{bib:2}N.\thinspace BOURBAKI, \emph{Algèbre commutative} (ibid)
(especially Chapter V, which deals with integrality, and Chapter VII, which
discusses Dedekind and factorial rings. Chapter II gives a very complete and
general discussion of rings of fractions and there is a good exposition of
valuation theory in Chapter VI).
\item[\textbf{[3]}] \phantomsection\label{bib:3}C.\thinspace CARATHÉODORY, \emph{Theory of functions},
Vol.~I (Chelsea, New York, 1954).
\item[\textbf{[4]}] \phantomsection\label{bib:4}S.\thinspace LANG, \emph{Algebraic structures}
(Addison-Wesley, 1967).
\item[\textbf{[5]}] \phantomsection\label{bib:5}S.\thinspace LANG, ``On quasi-algebraic closure'' (Ann.\
of Math., 55 (1962), 373--390).
\item[\textbf{[6]}] \phantomsection\label{bib:6}P.\thinspace SAMUEL, ``A propos du théorème des unités''
(Bull.\ Sci.\ Math., 90 (1966), 89--96).
\item[\textbf{[7]}] \phantomsection\label{bib:7}P.\thinspace SAMUEL, ``Anneaux factoriels'' (Publ.\ Soc.\
Math.\ São Paulo, 1964) (also see the article, ``Unique factorization,''
American Math.\ Monthly, 75 (1968), 945--952).
\item[\textbf{[8]}] \phantomsection\label{bib:8}G.\thinspace TERJANIAN, ``Sur une conjecture de
M.\thinspace Artin'' (C.R.\ Acad.\ Sci.\ Paris, (1966)).
\item[\textbf{[9]}] \phantomsection\label{bib:9}O.\thinspace ZARISKI and P.\thinspace SAMUEL,
\emph{Commutative algebra}, Vol.~I (Van Nostrand, Princeton, 1958)
(Chapter II discusses field theory, Chapter IV discusses Noetherian rings,
and Chapter V discusses integrality and Dedekind rings).
\end{list}
